% ---------------------------------------------------------------------------
% Appendix A: machine-verified mathematical core.
% Drop in after the main text, inside an \appendix environment.
% Requires amsmath, amssymb and hyperref.
% ---------------------------------------------------------------------------

\appendix
\section{Machine-verified mathematical core}
\label{app:lean}
\begingroup
\setlength{\emergencystretch}{3em}

The results of Sections~3--7 and 10 divide cleanly into two kinds of claim. One kind is
physical: that a speed--distance law of the form
$\varepsilon_n\tau_n\ge c_nd_{\rm phys,n}^2$ applies to a given class of
dynamics, that a chosen cost is legitimately charged against a cumulative
budget, and that a chosen physical metric controls operational
distinguishability. The other kind is mathematical: that the stated inequalities
entail the stated conclusions. Only the second kind can be machine-checked, and
this appendix reports that it has been.

\begin{quote}
The Lean formalization verifies the elementary mathematical implications
underlying the main no-go results, conditional on the stated hypotheses. It does
not formalize or validate the physical applicability of those hypotheses.
\end{quote}

Concretely, the mathematical core of the no-go argument was independently
formalized and machine-checked in Lean~4 with \texttt{Mathlib}. The development
contains no \texttt{sorry}, and \verb|#print axioms| reports that every result
below depends only on Lean's three standard axioms --- propositional
extensionality, quotient soundness, and choice --- so nothing rests on an
assumption added for convenience.

\subsection{What was formalized}

\paragraph{A.1 The material bound on the realization count.}
If $m_{\min}>0$ and, for every $n$, $m_{\min}N_n\le M_n$ and $M_n\le M_*$, then
\[
N_n\le\frac{M_*}{m_{\min}}\qquad\text{for every }n.
\]

\paragraph{A.2 The uniform rate bound.}
If $0\le N_n\le N_{\max}$, $0\le\nu_n\le\nu_{\max}$ and $0\le p_n\le1$, then
\[
\lambda_n=N_n\nu_np_n\le N_{\max}\nu_{\max}\qquad\text{for every }n.
\]
Only the upper bound $p_n\le1$ is used; no lower bound on $p_n$ is required.

\paragraph{A.3 The reciprocal-series obstruction (Theorem~3.1).}
If $0<\lambda_n\le\Lambda$ for every $n$, then
\[
\sum_n\frac1{\lambda_n}\text{ diverges}.
\]
This is formalized both in finite-prefix form and as
$\neg\,\mathrm{Summable}\,(n\mapsto1/\lambda_n)$. A combined declaration
derives the same conclusion directly from the material budget, material floor,
generation-rate ceiling, and $0\le p_n\le1$.

\paragraph{A.4 The weighted finite-action bound (Corollary~7.1).}
For nonnegative sequences $c_n,\varepsilon_n,\tau_n,d_n$ with
$c_nd_n^2\le\varepsilon_n\tau_n$ and summable $\varepsilon$ and $\tau$,
\[
\sum_n\sqrt{c_n}\,d_n
\le
\sqrt{\sum_n\varepsilon_n}\sqrt{\sum_n\tau_n}.
\]

\paragraph{A.5 The finite-action bound (Theorem~5.1).}
If in addition $c_n\ge c_*$ for a constant $c_*>0$, then $d$ is summable and
\[
\sum_n d_n
\le
\frac1{\sqrt{c_*}}
\sqrt{\sum_n\varepsilon_n}\sqrt{\sum_n\tau_n}.
\]
The exact variable-coefficient Lean declaration is given in the source as
\nolinkurl{finite_action_of_kinetic_floor}; the shorter
\texttt{finite\_action} is its constant-coefficient helper.

\paragraph{A.6 The fixed-resolution counting bound (Corollary~6.1).}
Under the hypotheses of A.5, for every finite set $s$ such that
$d_n\ge\delta$ for all $n\in s$,
\[
|s|\delta
\le
\frac1{\sqrt{c_*}}
\sqrt{\sum_n\varepsilon_n}\sqrt{\sum_n\tau_n},
\]
and consequently $\{n:d_n\ge\delta\}$ is finite. The exact manuscript-form
Lean declarations are listed in the source file as the counting theorem and
its corresponding finiteness theorem.

\paragraph{A.7 The finite-state operational bridge (Section~\ref{sec:metrics}).}
For a finite state space and a stochastic measurement channel $\Phi$,
\[
d_{\rm TV}(\Phi p,\Phi q)\le d_{\rm TV}(p,q).
\]
Hence any operational distance defined as the supremum of readout total
variation over a fixed admissible channel class satisfies
$d_{\rm op}\le d_{\rm TV}$. The Lean file
\nolinkurl{OperationalBridge.lean} verifies the $L^1$ contraction,
the total-variation data-processing inequality, sharpness for the identity
channel, transfer of a total-variation speed-limit hypothesis to
$d_{\rm op}$, and the fixed-resolution counting bound
\[
K_\delta^{\rm op}
\le
\frac1\delta\sqrt{\frac{\Sigma_*\mathcal N_*}{2}}.
\]
The stochastic-thermodynamic premise
$2d_{\rm TV}^2\le\Sigma\mathcal N$ is intentionally a hypothesis of the
Lean development rather than a theorem of it.

\subsection{What is not verified}

The following remain physical or modelling claims and are not certified by Lean:

\begin{enumerate}
\item the applicability of $\varepsilon_n\tau_n\ge c_nd_{\rm phys,n}^2$ to a
particular dynamics;
\item a universal identification or comparison of arbitrary physical metrics
with an operational organizational distance; the finite-state
total-variation/channel comparison of A.7 is machine-checked separately;
\item the claim that $\varepsilon_n$ is the correct cumulatively charged cost;
\item the stochastic-process theorem connecting the reciprocal-series criterion
to pure-birth explosion or non-explosion;
\item the physical additivity assumption $m_{\min}N_n\le M_n$.
\end{enumerate}

Thus this appendix does not show that the physical paper as a whole is formally
proved. It shows that the algebraic and analytic implications formalized here
are free of derivational gaps; the remaining burden lies in the explicitly
stated physical and stochastic-process premises.

\subsection{Reproduction}

The checked sources are
\nolinkurl{OrganizationalDepth.lean} and
\nolinkurl{OperationalBridge.lean}. They use Lean~4.33.0 and
\texttt{Mathlib} commit \texttt{db584cd} (10 August 2026). From the
verification directory, the check is reproduced with
\begin{verbatim}
lake update
lake exe cache get
lake build
\end{verbatim}
This exact three-command route is exercised by the repository CI. Both sources
compile with no errors and contain no \texttt{sorry}. The axiom audit for all
manuscript-facing declarations reports only
\texttt{propext}, \texttt{Classical.choice}, and \texttt{Quot.sound}.

The manuscript-facing declarations in \nolinkurl{OrganizationalDepth.lean}
cover stochastic-route closure, the weighted and uniform finite-action
bounds, and the fixed-resolution counting/finiteness results. The
declarations in \nolinkurl{OperationalBridge.lean} cover total-variation
data processing, identity-channel sharpness, transfer of a TV speed-limit
hypothesis to operational distance, and the operational fixed-resolution
counting bound.

\endgroup
